\documentclass[11pt,a4paper]{article}
\usepackage[UTF8]{ctex}
\usepackage{geometry}
\geometry{left=2.5cm,right=2.5cm,top=2.5cm,bottom=2.5cm}
\usepackage{amsmath,amssymb,amsthm}
\usepackage{hyperref}
\usepackage{booktabs}
\usepackage{enumitem}
\usepackage{xltabular}
\hypersetup{colorlinks=true,linkcolor=blue,urlcolor=blue}

\newtheorem{definition}{定义}[section]
\newtheorem{proposition}{命题}[section]
\newtheorem{remark}{注记}[section]

\title{SOAR 剧本优化与人机协同自适应切换（修正版）\\
——时间衰减、依赖关系、排队论与优雅降级}
\author{李龙胜}
\date{\today}

\begin{document}

\maketitle

\begin{center}
\textit{
SOAR 不再是简单的执行器，而是有限资源的优化调度者。\\
剧本执行有先后，响应收益随时间衰减。\\
剧本间存在依赖：先封禁还是先隔离？人工确认还是自动执行？\\
工单堆积时，等待时间越长，收益越低，分析师越疲劳。\\
最优调度顺序、最优人机切换阈值、最优等待时间——\\
恰好是三个边际均衡点的叠加。\\
当系统负载超过容量时，不是随机丢弃任务，\\
而是按价值排序、从低价值任务开始、逐步降级处理方式。\\
这是优雅降级——分析能力调度系统的普遍原则。\\
本文完成 SOAR 从执行者到优化者的完整升级。
}
\end{center}

\vspace{5mm}

\begin{abstract}
在已有的安全运营框架中，SOAR 作为裁定引擎的执行层，
负责执行自动剧本或派发人工工单。
然而，SOAR 本身的执行资源是有限的——
多个剧本竞争执行算力，人工研判队列存在排队延迟。
本文建立 SOAR 剧本调度的最优模型：
引入响应收益的时间衰减、剧本间的依赖与互斥关系、
人工研判队列的最优等待时间与分析师疲劳效应、
以及人机协同的自适应切换阈值。
SOAR 的最优调度顺序、人机切换阈值和人工等待时间，
均由贪心择优统一裁定，并与告警分析率 r* 实现联动调度。
在此基础上，提炼出“优雅降级”原则——
分析能力调度系统中资源紧缺时的普遍应对策略。
\end{abstract}

\tableofcontents

\section{引言}

在现有的四层分析能力分配框架中，
裁定引擎输出裁定结果——忽略、自动分析、人工分析。
SOAR 作为执行层，负责执行自动分析对应的剧本，
或将人工分析派发为工单。

然而，SOAR 本身也面临资源分配问题。
当裁定引擎同时输出多个“自动分析”指令时，
SOAR 必须在这些剧本之间调度执行顺序。
同时执行的剧本数量受并发限制和 API 速率限制约束，
剧本执行本身消耗计算资源，
与告警分析共享同一块分析能力预算。

此外，人工研判队列同样面临排队延迟。
当工单积压时，新工单需要等待，
等待时间越长，响应收益衰减越多，
分析师的判断准确率也随疲劳而下降。
SOAR 需要动态切换自动执行与人工研判的比例，
在响应速度与误报风险之间找到最优平衡。

本文将 SOAR 从执行者升级为优化者，
完成其与已有框架的完整对接。

\section{最优剧本调度}

\subsection{响应收益的时间衰减}

设当前有 \(N\) 个待执行的 SOAR 剧本。
每个剧本 \(j\) 关联一个响应收益 \(V_j\)——
成功执行该剧本能避免的潜在安全损失。
响应收益不是恒定的，它随时间 \(t\) 衰减：
\[
V_j(t) = V_j(0) \cdot e^{-\lambda_j t},
\]
其中 \(\lambda_j\) 是衰减速率，度量攻击影响的紧迫度。
正在进行的实时攻击，\(\lambda_j\) 很大，收益在秒级衰减。
已完成但尚未清理的后门，\(\lambda_j\) 较小，收益在小时级衰减。

\subsection{调度优先级}

设剧本 \(j\) 的执行成本为 \(c_j\)。
在忽略时间衰减时，最优调度按收益成本比 \(V_j / c_j\) 降序执行。
引入时间衰减后，调度还必须考虑时间窗口。
设剧本 \(j\) 已等待时间为 \(w_j\)，
当前实际收益为 \(V_j(w_j)\)，
调度优先级修正为：
\[
\text{Priority}_j = \frac{V_j(w_j)}{c_j}.
\]

最优调度规则：
\begin{enumerate}[label=(\arabic*)]
    \item 优先执行时间窗口即将关闭的剧本——
          即使其绝对收益较低，因其时间紧迫度更高。
    \item 当多个剧本的时间窗口都充裕时，按绝对收益成本比降序执行。
    \item 当系统负载过高、多个剧本的时间窗口都可能错过时，
          必须舍弃那些即使立即执行也已无法挽回大部分损失的剧本——
          集中资源保护尚有时间窗口的剧本。
\end{enumerate}

这是贪心择优在时间维度上的扩展——
每步选择当前边际收益最高的剧本，
同时计入时间衰减效应。

\textbf{时间衰减与依赖关系的联合优化}：
时间衰减和依赖关系在实时系统中同时存在。
一个剧本时间窗口很紧迫，但其前置依赖剧本尚未完成；
另一个剧本时间窗口稍宽松，但已可立即执行。
贪心策略必须同时权衡这两个约束——
优先执行那些当前可执行且时间窗口最紧迫的剧本。
时间衰减与依赖关系的联合优化
贪心策略提供的是近似最优解。
精确最优性的证明或反例尚待形式化分析，
当前采用贪心启发式。

\section{剧本间的依赖关系}

之前假设剧本之间相互独立。现实中剧本间存在复杂依赖：

\textbf{依赖关系}：
剧本 B 必须在剧本 A 完成后才能执行。
例如，必须先分析文件获取哈希，
才能根据哈希封禁相关终端。

\textbf{互斥关系}：
两个剧本不能同时执行。
例如，不能同时对同一终端执行隔离和深度扫描。

\textbf{级联效应}：
执行剧本 A 后，可能自动触发剧本 B 和剧本 C。
例如，隔离一台终端后，需要重新评估该终端所在网段的威胁等级，
可能触发新的扫描。

这些依赖关系使 SOAR 剧本调度
从简单的排序问题升级为有向无环图的拓扑排序问题。
最优调度需要在尊重依赖关系和互斥约束的前提下，
最大化总响应收益。

实时系统中，当剧本依赖图规模较小时，
可以使用精确求解；规模较大时，
应使用启发式近似求解以保证实时性。
具体算法选择在工程实现中根据实际部署场景确定。

\section{人工研判队列的最优等待时间}

人工研判队列面临排队论的经典问题。
当工单到达速率超过分析师服务速率时，
队列无限增长，每个工单的等待时间趋于无穷，
响应收益趋于零。

当告警洪流到来时，人工研判队列迅速饱和，
大量需要人工确认的剧本卡在队列中，
等待分析师确认，
而攻击方在等待期间已经达成目标。

\subsection{人工研判队列的贪心择优调度}

工单进入队列前，先按预期响应收益和等待成本进行筛选。
低价值工单直接取消人工确认环节，
降级为自动执行或归档。
高价值工单保留人工确认，
但设置最大等待时间 \(T_{\max}\)——
超过此时间后自动降级为自动执行，
避免响应收益因无限等待而完全消失。

最大等待时间的最优设定 \(T_{\max}^*\)
是边际等待收益等于边际误报风险的均衡点。
等待太久收益衰减殆尽，等待不够误报风险升高。
这是够用就行原则在人工研判队列维度的直接应用。

\subsection{差异化最大等待时间}

在实际实现中，\(T_{\max}\) 不是单一值，
而是告警类别的函数 \(T_{\max}(i)\)。
高紧迫度类别的 \(T_{\max}\) 较小，
低紧迫度类别的 \(T_{\max}\) 较大。
各 \(T_{\max}(i)\) 的最优值
由各类别的响应收益衰减速率 \(\lambda_i\)
和误报风险 \(\alpha_i\) 共同决定——
与已有框架中告警分析率 r* 的类别差异化裁定完全同构。

\subsection{分析师疲劳效应}

当前模型假设分析师的服务质量恒定。
在真实场景中，分析师在连续研判高复杂度工单后，
判断准确率会下降。
这意味着人工研判的服务质量
随队列长度和研判强度衰减。
排在队列后面的工单即使等到分析师，
也可能被疲劳的分析师误判——
误判风险远高于正常状态。
这一效应当前未被显式建模，
需心理学和工效学实证支持。

\section{人机协同自适应切换阈值}

SOAR 剧本中，部分必须人工确认后执行，
部分可以自动执行，部分处于灰色地带。
人机协同的核心是切换阈值的动态裁定。

最优切换阈值不是固定的，
它取决于当前分析能力的饱和度、
攻击方探测强度和剧本本身的误操作风险：
\begin{itemize}
    \item 分析能力充裕时，阈值偏向人工研判——
          更谨慎，允许更多灰色地带事件进入人工队列。
    \item 告警洪流来临时，阈值偏向自动执行——
          更快速，只有最高风险的操作才保留人工确认环节。
    \item 对手行为分析检测到探测强度上升时，
          响应速度优先，阈值自动偏向自动执行。
    \item 攻击方蛰伏时，阈值自动恢复为偏人工研判，
          节省自动化资源。
\end{itemize}

\textbf{误操作风险的差异化处理}：
人机切换阈值是告警类别 \(i\) 的函数——
不仅取决于当前分析能力饱和度 \(s(t)\)
和攻击方探测强度，
还取决于剧本 \(i\) 的误操作风险 \(M_i\)。
\(M_i\) 越高的剧本，其自动执行阈值越严格。
即使在告警洪流中，
高误操作风险的剧本也应保留人工确认环节。
\(M_i\) 的评估来自台账中该类告警的 \(A_i\)
和误操作影响范围。

最优切换阈值由边际响应速度收益
等于边际误报风险的均衡条件决定——
这是够用就行原则在 SOAR 维度的直接应用。

\section{SOAR 与告警分析率 r* 的联动调度}

裁定引擎的最优分析率 r* 决定“哪些告警值得被处理”，
SOAR 调度决定“这些值得被处理的告警中，
先处理哪个、自动还是人工”。
两者之间存在天然的联动关系。

当总分析能力预算紧张时：
\begin{itemize}
    \item 裁定引擎降低低价值告警的 r*，
          从源头减少进入 SOAR 的告警量。
    \item SOAR 自动调高自动执行的比例，
          减少人工排队。
    \item 人工研判队列自动缩短最大等待时间，
          让积压工单降级为自动执行或归档。
\end{itemize}

三层协同由贪心择优统一裁定——
r* 决定告警是否进入 SOAR 队列，
调度优先级决定进入后何时执行，
人机切换阈值决定执行方式。
三者在分析能力池的统一约束下达到边际均衡。

\section{优雅降级——分析能力调度系统的普遍原则}

在告警分析中，降级意味着减少投入——
从人工分析降为自动分析，从自动分析降为忽略。
在 SOAR 中，降级意味着减少等待——
从“等人来确认”降为“机器直接做”。
在蜜罐运维中，降级意味着减少真实性投入——
从高仿真蜜罐降为低交互蜜罐。
在零知识审计中，降级意味着扩大聚合批次——
从逐条生成证明降为批量聚合证明。
在 UEBA 中，降级意味着提高筛选阈值——
从捕获所有偏离事件降为仅捕获高偏离事件。

这些降级操作共享同一个逻辑结构：
当系统负载超过容量时，不是随机丢弃任务，
而是按价值排序、从低价值任务开始、
逐步降级处理方式。
这是分析能力调度系统的普遍原则——
优雅降级。

优雅降级的核心是价值排序。
系统必须知道哪些任务是高价值的、哪些是低价值的，
才能在降级时从低价值开始。
这个价值排序恰好来自台账中的 A 值——
漏报损失越高的告警类别，其处理绝不能被降级。
够用就行原则在优雅降级中体现为：
降级的边际节省必须大于边际损失，
当两者相等时，系统在当前负载下达到了最优降级水平。

\section{诚实标注}

\begin{center}
\begin{xltabular}{\textwidth}{c|X|c}
\toprule
\textbf{命题} & \textbf{状态} & \textbf{层级} \\
\midrule
SOAR 剧本最优调度（含时间衰减） &
贪心择优在时间维度上的扩展，数学结构与已有模型完全同构 &
II \\
时间衰减与依赖关系的联合优化 &
贪心策略提供近似最优解，
精确最优性的证明或反例尚待形式化分析 &
II（当前采用贪心启发式） \\
剧本间依赖关系的处理 &
识别为有向无环图拓扑排序问题，
约束优化框架已明确 &
II（精确求解算法待工程实现中选定） \\
人工研判队列最优等待时间 &
排队论框架下的贪心择优调度，均衡条件已明确，
差异化 \(T_{\max}(i)\) 与已有框架同构 &
II \\
分析师疲劳效应 &
当前模型假设服务质量恒定，
疲劳导致的质量衰减未被显式建模 &
III（需心理学和工效学实证支持） \\
人机协同自适应切换阈值 &
动态阈值与分析能力饱和度、攻击方探测强度和误操作风险挂钩，
逻辑清晰 &
II \\
误操作风险的差异化处理 &
\(M_i\) 的维度已引入，
高误操作风险剧本在洪流中也保留人工确认 &
II（\(M_i\) 的定量评估待台账标注） \\
SOAR 与告警分析率 r* 的联动调度 &
三层协同由贪心择优统一裁定，逻辑自洽 &
II \\
优雅降级原则 &
识别为告警分析、SOAR、蜜罐、审计、UEBA
五个维度的共享逻辑结构 &
II（定性规律已确立，形式化表达待展开） \\
\bottomrule
\end{xltabular}
\end{center}

\section{结语}

SOAR 剧本优化完成了 SOAR 从执行者到优化者的升级。
最优调度顺序是时间衰减、依赖关系和响应收益成本的联合最优解。
人工研判队列的最优等待时间
由各告警类别的响应收益衰减速率和误报风险差异化决定。
人机协同的最优切换阈值
随分析能力饱和度、攻击方探测强度和剧本误操作风险自适应调节。
SOAR 与告警分析率 r* 的联动调度，
由贪心择优在分析能力池的统一约束下裁定。

优雅降级原则的提炼揭示了
告警分析、SOAR、蜜罐、审计和 UEBA
五个看似独立的模块
共享同一个深层逻辑结构——
当资源紧缺时，按价值排序、从低价值任务开始、逐步降级处理方式。
这是够用就行原则在资源紧缺场景中的最通用表达。

\vspace{5mm}
\noindent\textbf{文档信息}：
\begin{itemize}
    \item 作者：李龙胜
    \item 版本：修正版（整合外部审查反馈）
    \item 许可：CC BY 4.0
    \item 前置依赖：四卷体系、安全参数自适应调节文档、产品蓝图修正版
\end{itemize}

\end{document}